\section{Meta-RL 的理论视角}

\subsection{Meta-RL 作为 POMDP}

\begin{importantbox}{Meta-RL 的 POMDP 视角}
可以将 meta-RL 视为在一个更大的 POMDP 中做 RL：
\begin{itemize}
    \item 隐状态包含任务身份
    \item 观测是当前状态和历史经验
    \item 最优策略需要在交互过程中推断任务身份（贝叶斯推断）
\end{itemize}
从这个视角看，context-based 方法本质上是在近似这个 POMDP 的最优策略。
\end{importantbox}

\subsection{后验采样与 Thompson Sampling}

在理想情况下，meta-RL 的最优行为应该类似于 Thompson Sampling：
\begin{enumerate}
    \item 维护对任务身份的后验信念 $p(\mathcal{T} | \text{history})$
    \item 从后验中采样一个任务假设
    \item 按照该假设下的最优策略行动
    \item 收集新数据后更新后验
\end{enumerate}

\subsection{本章小结}

从理论角度看，Meta-RL 可以统一理解为在任务不确定的 POMDP 中做最优决策，这为理解和改进 meta-RL 方法提供了原则性的框架。

